\section{Conclusion}
\label{sec:conclusion}

We presented a dual-mode multimodal Next Best Action (NBA) recommendation system that extracts features from
text, visual, and behavioral multimedia data to unify active search and
passive personalization within a single architecture.
The system integrates multimodal feature extraction via BGE-M3 and CLIP
with co-purchase graph embeddings and DIF-SASRec sequential personalization,
evaluated end-to-end against multiple baseline approaches.

By conditioning DIF-SASRec on dense BGE-M3 content embeddings rather than
item IDs, the system generalises to the full 3.08M-item catalog without
retraining.
Pipeline~A (behavioral graph embedding + BGE-M3 ranking) provides strong
content-aware retrieval, while DIF-SASRec delivers complementary
sequential personalization via online adaptation.
Complementarity analysis shows Pipeline~A recovers 88.3\% of DIF-SASRec
misses (Table~\ref{tab:complementarity}), and the union system achieves HR@10 of 0.9736
(+124\% over the content-only baseline), with warm-cache P50 latency of
34\,ms and safe concurrent inference via the model pool.

Future work includes dynamic graph embedding retraining from accumulated
interaction logs, expansion of multilingual keyword search beyond
the current inverted index, cross-domain transfer of pretrained
DIF-SASRec weights, and extension of the NBA framing to domains
beyond book discovery.
These results demonstrate that fusing heterogeneous retrieval signals within a
unified, continuously updated user model is a practical and scalable path toward
NBA systems capable of serving both explicit and exploratory information needs
across large-scale multimedia catalogs.
